\documentclass[11pt]{book}
\usepackage[paperwidth=145mm,paperheight=200mm,top=18mm,bottom=20mm,inner=20mm,outer=16mm]{geometry}
\usepackage{brumfiel}
% figures17.tex -- Appendix figures redrawn in TikZ.
%
% The Appendix (book pp.275-282) contains exactly ONE figure: the unnumbered
% and uncaptioned diagram on book p.277 that accompanies Postulate 16. It
% shows two three-ray configurations, unprimed and primed, drawn congruent.
% (Pages 275, 276, 278-281 are text only; 282 is blank. Verified page by page
% against sources/Geometry.pdf pp.289-296.)
%
% GEOMETRY MEASURED off sources/Geometry.pdf p.291 rasterised at 600 dpi.
% Method: threshold the ink, separate the heavy rays from the light arcs by
% binary erosion (rays survive a 5x5 erosion, arcs do not), take the vertex as
% the left end of the horizontal ray, then fit ray directions from the eroded
% ray cores only. Reading the whole ink mask instead biases every direction
% toward the arcs that run tangent to it -- that is what produced the earlier
% 45.0/117.0 figures. Raw and de-skewed (photo tilt removed by subtracting the
% measured tilt of the horizontal ray) results:
%
%     ray          raw unprimed  raw primed   de-skewed mean
%     horizontal       +0.48        +0.67          0 (datum)
%     diagonal         44.38        43.87        43.55
%     upper-left      116.33       116.13       115.66
%
% Drawn at 0 / 43.6 / 115.7 deg. Ray lengths and arc radii as fractions of the
% horizontal ray (mean of the two sub-figures):
%
%     horizontal ray 1.000   arc A 0.446   (drawn 3.30 / 1.47)
%     diagonal ray   0.994   arc C 0.699   (drawn 3.28 / 2.31)
%     upper-left ray 0.939   arc B 0.774   (drawn 3.10 / 2.55)
%
% The three arc radii are DISTINCT and ordered A < C < B, and each label sits
% on its own arc in a gap in that arc. Arc spans read off the source by
% scanning each arc's circle for ink (rays and label glyphs masked out):
%
%     arc A   4..60  and  78..114   gap 60..78, label at  69 deg
%     arc B  47..77  and  87..114   gap 77..87, label at  81.5 deg
%     arc C   4..16  and  30..42    gap 16..30, label at  23 deg
%
% The arrowhead ends below are the source's exactly; the GAP ends are opened a
% few degrees wider (A 52..86, B 71.5..91.5, C 12..34) because our labels are
% set with more padding than the book's tight hand lettering -- at the source's
% own gap widths check_labels.py reports the glyph boxes touching the arcs.
%
% Arrowheads therefore land at 114 (upper-left ray), 47/42 (diagonal ray) and
% 4 (horizontal ray) -- always just short of the ray they point at, and the two
% hits on the horizontal ray (arc A at 1.47, arc C at 2.31) and the two on the
% diagonal (arc B at 2.55, arc C at 2.31) are at different radii, as in the book.
%
% LINE WEIGHT: the book draws the three rays roughly 1.75x the weight of the
% annotation arcs (measured 13.0 px vs 8.5 px of ink at 600 dpi on the same
% photo). The rays keep the house `given' weight so the Appendix matches the
% rest of the book; the arcs are dropped to 0.42pt to restore the contrast.
%
% What the postulate rests on: the diagonal ray is INTERIOR to angle A, so
% angle A = angle B + angle C. See tools/constraints/ch17.py.

% ------------------------------------------------ Appendix figure (p.277)
\newcommand{\FIGAPPONE}{\begin{bkfigure}[0.52\linewidth]
% --- unprimed configuration
\begin{tikzpicture}[scale=0.72,
    mk/.style={fig, line width=0.42pt},          % light annotation arcs
    alab/.style={vlab, outer sep=1.5pt}]         % label sits in its arc's gap
  \coordinate (O) at (0,0);
  \draw[given] (O) -- (0:3.30);        % horizontal ray
  \draw[given] (O) -- (43.6:3.28);     % diagonal ray, interior to angle A
  \draw[given] (O) -- (115.7:3.10);    % upper-left ray
  % angle A : upper-left ray <-> horizontal ray  (the whole angle)
  \draw[mk,->] (86:1.47) arc[start angle=86, end angle=114, radius=1.47];
  \draw[mk,->] (52:1.47) arc[start angle=52, end angle=4,   radius=1.47];
  \node[alab] at (69:1.47) {$A$};
  % angle B : upper-left ray <-> diagonal ray
  \draw[mk,->] (91.5:2.55) arc[start angle=91.5, end angle=114, radius=2.55];
  \draw[mk,->] (71.5:2.55) arc[start angle=71.5, end angle=47,  radius=2.55];
  \node[alab] at (81.5:2.55) {$B$};
  % angle C : diagonal ray <-> horizontal ray
  \draw[mk,->] (34:2.31) arc[start angle=34, end angle=42, radius=2.31];
  \draw[mk,->] (12:2.31) arc[start angle=12, end angle=4,  radius=2.31];
  \node[alab] at (23:2.31) {$C$};
\end{tikzpicture}

\vspace{4mm}

% --- primed configuration (drawn congruent to the unprimed one, as the book
%     draws it; the postulate itself only supposes congruence of two pairs)
\begin{tikzpicture}[scale=0.72,
    mk/.style={fig, line width=0.42pt},
    alab/.style={vlab, outer sep=1.5pt}]
  \coordinate (O) at (0,0);
  \draw[given] (O) -- (0:3.30);
  \draw[given] (O) -- (43.6:3.28);
  \draw[given] (O) -- (115.7:3.10);
  \draw[mk,->] (86:1.47) arc[start angle=86, end angle=114, radius=1.47];
  \draw[mk,->] (52:1.47) arc[start angle=52, end angle=4,   radius=1.47];
  \node[alab] at (69:1.47) {$A'$};
  \draw[mk,->] (91.5:2.55) arc[start angle=91.5, end angle=114, radius=2.55];
  \draw[mk,->] (71.5:2.55) arc[start angle=71.5, end angle=47,  radius=2.55];
  \node[alab] at (81.5:2.55) {$B'$};
  \draw[mk,->] (34:2.31) arc[start angle=34, end angle=42, radius=2.31];
  \draw[mk,->] (12:2.31) arc[start angle=12, end angle=4,  radius=2.31];
  \node[alab] at (23:2.31) {$C'$};
\end{tikzpicture}
\end{bkfigure}}


% ---- local macros (hoist into book preamble) ----
% The Appendix has no numbered "N--M" sections. Its three major heads are
% set centred, bold and small, not flush left like a chapter's \section*.
% \APPhead also files the TOC entry the assembled book needs.
\newcommand{\APPhead}[2]{%
  \par\bigskip\begin{center}\bfseries #1\end{center}%
  \addcontentsline{toc}{section}{#2}\par\nopagebreak\medskip}
% Bold centred subhead grouping a run of postulates.
\newcommand{\APPsub}[1]{%
  \par\medskip\begin{center}\bfseries\small #1\end{center}\par\nopagebreak}
% Continuously-numbered reference list, with (a)/(b)/(c) subparts.
\newlist{applist}{enumerate}{2}
\setlist[applist,1]{label=\arabic*., leftmargin=1.7em, labelsep=0.4em,
                    itemsep=2.5pt, topsep=4pt, parsep=0pt}
\setlist[applist,2]{label=(\alph*), leftmargin=1.9em, labelsep=0.4em,
                    itemsep=1.5pt, topsep=2.5pt, parsep=0pt}
% ---- end local macros ----

\begin{document}
\setcounter{chapter}{17}
% The Appendix text is verbatim, so line breaking must bend, not the wording:
% give TeX a little extra stretch rather than reword to clear overfull boxes.
\emergencystretch=2em

% The book sets this opener flush RIGHT: "APPENDIX" occupies the slot that
% "CHAPTER 10" occupies on a numbered chapter opener (cf. book p.180).
\chapter*{\raggedleft\Huge APPENDIX}
\addcontentsline{toc}{chapter}{Appendix}
\markboth{APPENDIX}{}
\vspace{-1em}\noindent\rule{\linewidth}{0.6pt}\vspace{1em}
\figurekey

This appendix lists all the postulates for plane geometry, along with
some of the definitions and theorems. In addition to being convenient
for reference purposes, the appendix merits study, because reading it
over and thinking about the order of the theorems will help you fix the
structure of Euclidean geometry in your mind.

%========================================================= the postulates
\APPhead{THE POSTULATES}{The Postulates}

There are three groups of postulates, and three postulates listed singly.
They are:

\begin{center}
\begin{minipage}{0.86\linewidth}
Postulates of incidence (3 postulates)\\
Postulates of order or betweenness (5 postulates)\\
Postulates of congruence (7 postulates)\\
The postulate of Archimedes\\
The completeness postulate\\
The parallel postulate.
\end{minipage}
\end{center}

\APPsub{The Incidence Postulates}

\begin{applist}
\item There are at least three points not all on a line.
\item For any two different points, there is exactly one line containing
      these points.
\item Every line contains at least two points.
\end{applist}

\APPsub{The Betweenness Postulates}

\begin{applist}[resume]
\item If $B$ is between $A$ and $C$, then $A$, $B$, and $C$ are three
      different, or distinct, points on a line.
\item For every three points on a line, exactly one of them is between the
      other two.
\item Any four points on a line may be named $A_1$, $A_2$, $A_3$, $A_4$, so
      that the only betweenness relations are the same as the order of the
      subscript numbers. That is, the betweenness relations are: $A_2$
      between $A_1$ and $A_3$; $A_2$ between $A_1$ and $A_4$; etc.
\item If $A$ and $B$ are two points, then there is at least one point $C$
      such that $B$ is between $A$ and $C$, and at least one point $D$ such
      that $D$ is between $A$ and $B$.
\item Every line separates the plane. By this we mean that all points of
      the plane not on the line are divided into two sets, called the two
      sides of the line, having the following properties:
      \begin{applist}
      \item If $P$ and $Q$ belong to one of these sets, then no point of
            $l$ is between $P$ and $Q$. We say then that $P$ and $Q$ are on
            the same side of $l$.
      \item If $P$ and $Q$ are in different sets, then there is a point on
            $l$ which is between $P$ and $Q$.
      \end{applist}
\end{applist}

\APPsub{Linear Congruence Postulates}

\begin{applist}[resume]
\item Given points $A$ and $B$ on a line $l$ and a point $A'$ on a line
      $l'$, then on a given side of $A'$ on $l'$ there is exactly one point
      $B'$ such that $AB$ is congruent to $A'B'$. To indicate that $AB$ is
      congruent to $A'B'$, we write ``$AB \cong A'B'$.''

      In congruence the order of the points does not matter, that is, if
      $AB \cong A'B'$, then also $AB \cong B'A'$, $BA \cong A'B'$, and
      $BA \cong B'A'$.
\item Congruence also satisfies the following:
      \begin{applist}
      \item $AB \cong AB$.
      \item If $AB \cong A'B'$, then $A'B' \cong AB$.
      \item If $AB \cong A'B'$ and $A'B' \cong A''B''$, then
            $AB \cong A''B''$.
      \end{applist}
\item (This postulate tells how to ``add'' and ``subtract'' segments.)
      Suppose that $B$ is between $A$ and $C$ on a line $l$, and that $B'$
      is between $A'$ and $C'$ on a line $l'$.
      \begin{applist}
      \item If $AB \cong A'B'$ and $BC \cong B'C'$, then $AC \cong A'C'$.
      \item If $AC \cong A'C'$ and $BC \cong B'C'$, then $AB \cong A'B'$.
      \end{applist}
\end{applist}

\APPsub{Archimedes' Postulate}

\begin{applist}[resume]
\item If $AB$ is any segment on a line $l$ and $CD$ any other segment, then
      there is a finite number $n$ of points $A_1, A_2, \ldots, A_n$ on $l$
      such that the points $A, A_1, A_2, \ldots, A_n$ are arranged in this
      order, and all the segments $AA_1, A_1A_2, \ldots, A_{n-1}A_n$ are
      congruent to $CD$, and either $B = A_n$ or $B$ lies between $A_{n-1}$
      and $A_n$.
\end{applist}

\APPsub{The Completeness Postulate}

\begin{applist}[resume]
\item For every line $l$ and any given point $A$ on $l$, and for any
      positive real number $x$, there is, on a given side of $A$, a point
      $B$ such that $\sg{AB} = x$.
\end{applist}

\APPsub{Congruence Postulates for Angles}

\begin{applist}[resume]
\item If $\ang{A}$ is not a straight angle and if $r'$ is a ray from $A'$ on
      a line $l'$, then on a given side of $l'$ there is exactly one ray
      $s'$ from $A'$ such that the angle $A'$, with sides $r'$ and $s'$, is
      congruent to angle $A$. We write this as ``$\ang{A'} \cong \ang{A}$.''
      If $\ang{A}$ is a straight angle, then the straight angle at $A'$ with
      one side $r'$ is the only angle with one side $r'$ congruent to angle
      $A$.
\item ~
      \begin{applist}
      \item $\ang{A} \cong \ang{A}$.
      \item If $\ang{A} \cong \ang{B}$, then $\ang{B} \cong \ang{A}$.
      \item If $\ang{A} \cong \ang{B}$ and $\ang{B} \cong \ang{C}$, then
            $\ang{A} \cong \ang{C}$.
      \end{applist}
\item If in the figure any two of the pairs $\ang{A}$ and $\ang{A'}$,
      $\ang{B}$ and $\ang{B'}$, $\ang{C}$ and $\ang{C'}$ are congruent, then
      the third pair of angles are congruent. In other words:
      \begin{applist}
      \item If $\ang{A} \cong \ang{A'}$ and $\ang{B} \cong \ang{B'}$, then
            $\ang{C} \cong \ang{C'}$.
      \item If $\ang{B} \cong \ang{B'}$ and $\ang{C} \cong \ang{C'}$, then
            $\ang{A} \cong \ang{A'}$.
      \item If $\ang{C} \cong \ang{C'}$ and $\ang{A} \cong \ang{A'}$, then
            $\ang{B} \cong \ang{B'}$.
      \end{applist}
      The angle $A$ may be a straight angle.
\end{applist}

\FIGAPPONE

\begin{applist}[resume]
\item If for $\tri{ABC}$ and $\tri{A'B'C'}$, $AB \cong A'B'$,
      $AC \cong A'C'$, and $\ang{A} \cong \ang{A'}$, then
      $\ang{B} \cong \ang{B'}$ and $\ang{C} \cong \ang{C'}$.
\end{applist}

\APPsub{The Parallel Postulate}

\begin{applist}[resume]
\item Through a point not on a line there is no more than one line parallel
      to the given line.
\end{applist}

%========================================================= definitions
\APPhead{SELECTED DEFINITIONS}{Selected Definitions}

Although every definition has a purpose, certain ones are more necessary
than others, and most of these essential definitions are listed here. In
some cases an abbreviated form of the definition is given when no
misunderstanding can arise.

\begin{applist}
\item Each set of points in the plane is called a \emph{geometric figure}.
\item If $A$ and $B$ are any two points on a line, then the set of points
      consisting of $A$ and $B$ and all the points between $A$ and $B$ will
      be called the \emph{line segment} (or \emph{segment}) determined by
      $A$ and $B$. The segment will be indicated by ``$AB$.'' The points
      $A$ and $B$ are called the \emph{ends} of the segment.
\item If $A$, $B$, and $C$ are three points not all on a line, the set of
      points of the segments $AB$, $BC$, $AC$ is called a \emph{triangle}.
      We call the points $A$, $B$, and $C$ \emph{vertices} and the segments
      $AB$, $AC$, $BC$ \emph{sides} of the triangle. To indicate ``triangle
      $ABC$,'' we write ``$\tri{ABC}$.''
\item If $O$ is a point of a line $l$ and $A$ and $B$ are two other points
      of $l$, then we say that $A$ and $B$ are on the \emph{same side of
      $O$} if $O$ is not between $A$ and $B$. We say that $A$ and $B$ are on
      \emph{opposite sides of $O$} if $O$ is between $A$ and $B$.
\item A \emph{ray} from $O$, or a ray with end $O$, is a set of
      points consisting of $O$ and all the points on one side of $O$ of a
      line $l$ through $O$.
\item The set of all points on two rays from a point $O$ is called an
      \emph{angle}. The point $O$ is called the \emph{vertex} of the angle,
      and the two rays are called the \emph{sides} of the angle. If the two
      rays are on the same line, the angle is called a \emph{straight
      angle}.
\item Consider an angle, not a straight angle, with vertex $O$ and whose
      sides are the rays $r_1$ and $r_2$ on the lines $l_1$ and $l_2$,
      respectively. The \emph{interior of the angle} is the set of all
      points on the same side of $l_1$ as $r_2$ and on the same side of
      $l_2$ as $r_1$. The \emph{exterior of the angle} is the set of all
      points in the plane which are not points of the angle and are not in
      the interior of the angle.
\item Two angles are \emph{adjacent} to each other if they have a common
      side and vertex and have no common interior points.
\item Let $AB$ be any segment and $A'B'$ any other segment on a line $l'$.
      On the same side of $A'$ as $B'$, let $B''$ be the point such that
      $AB \cong A'B''$. If $B'$ is between $A'$ and $B''$, we say that
      $AB > A'B'$. If $B''$ is between $A'$ and $B'$, we say that
      $AB < A'B'$.
\item For the definition of \emph{length of a segment} read Chapter 5,
      Section 4.
\item A \emph{right} angle is one which is congruent to its supplement.
\item Two triangles are said to be \emph{congruent} if they can be labeled
      $\tri{ABC}$ and $\tri{A'B'C'}$ such that $AB \cong A'B'$,
      $AC \cong A'C'$, $BC \cong B'C'$, $\ang{A} \cong \ang{A'}$,
      $\ang{B} \cong \ang{B'}$, and $\ang{C} \cong \ang{C'}$. To indicate
      that ``triangle $ABC$ is congruent to triangle $A'B'C'$,'' we write
      ``$\tri{ABC} \cong \tri{A'B'C'}$.''
\item Let $A$ and $B$ be any angles, not straight angles. Let $\ang{B'}$ be
      an angle congruent to $\ang{B}$ such that one side of $\ang{B'}$ is
      the same as the first side of $\ang{A}$, and the other side of
      $\ang{B'}$ and the second side of $\ang{A}$ lie on the same side of
      the first side of $\ang{A}$. If the second side of $\ang{B'}$ is
      interior to $\ang{A}$, we say $\ang{B} < \ang{A}$. If the second side
      of $\ang{B'}$ is exterior to $\ang{A}$, we say $\ang{B} > \ang{A}$. If
      $\ang{A}$ is a straight angle and $\ang{B}$ is not a straight angle,
      we say $\ang{A} > \ang{B}$ and $\ang{B} < \ang{A}$.
\item Let $P_1, P_2, \ldots, P_n$ be any set of $n$ points. The set of
      points on all the segments $P_1P_2, P_2P_3, \ldots, P_{n-1}P_n$ is
      called a \emph{polygonal path}. The path is said to join, or connect,
      the endpoints $P_1$ and $P_n$. The points $P_1, \ldots, P_n$ are
      called the \emph{vertices}. In particular, if $P_1 = P_n$, the
      polygonal path is called a \emph{polygon} and designated as polygon
      $P_1P_2 \ldots P_{n-1}$.

      A polygon is \emph{convex} if for every two points on the polygon each
      point between them is either on the polygon or in the interior. A
      convex, equiangular, and equilateral polygon is called \emph{regular}.
\item Two lines are \emph{parallel} if they do not intersect.
\item If the ray $OB$ is in the interior of $\ang{AOC}$ (or $\ang{AOC}$ is
      straight), then $\ang{AOC}$ is the \emph{sum} of $\ang{AOB}$ and
      $\ang{BOC}$.
\item If $AB$ and $CD$ are segments, the \emph{ratio} of $AB$ to $CD$ is the
      number $\sg{AB} \div \sg{CD}$. The statement that the segments $AB$
      and $CD$ are \emph{proportional} to the segments $A'B'$ and $C'D'$
      means that $\sg{AB}/\sg{A'B'} = \sg{CD}/\sg{C'D'}$. We say that the
      sides of $\tri{ABC}$ are proportional to the sides of $\tri{A'B'C'}$
      if $\sg{AB}/\sg{A'B'} = \sg{BC}/\sg{B'C'} = \sg{CA}/\sg{C'A'}$.
\item A \emph{circle} is determined by a point $O$, called the center, and a
      positive number $r$. The circle is the set of all points $A$ of the
      plane such that the distance $\sg{OA} = r$. If $A$ is a point of the
      circle, then the segment $OA$ is called a \emph{radius} (plural,
      \emph{radii}). If $A$ and $B$ are on the circle and the segment $AB$
      contains $O$, then $AB$ is called a \emph{diameter}.
\item If a circle has center $O$ and radius $r$, the \emph{interior of the
      circle} is the set of all points $B$ such that $\sg{OB} < r$. The
      \emph{exterior of the circle} is the set of all points $C$ such that
      $\sg{OC} > r$. A \emph{chord} is a segment determined by two points
      $E$ and $F$ on the circle. The line through $E$ and $F$ is a
      \emph{secant}. A \emph{central angle} of a circle is an angle whose
      vertex is at the center of the circle. An \emph{inscribed angle} is an
      angle with vertex on the circle and a second point of the circle on
      each of the two sides of the angle.
\item Suppose that for each positive integer $n$ there is associated in some
      way, a number $x_n$. The numbers $x_1, x_2, \ldots, x_n$ then form an
      \emph{infinite sequence}. The number $x_n$ is called the $n$th term of
      the sequence. We designate the sequence by the symbol $\{x_n\}$.
\item The number $L$ is called the \emph{limit of the sequence} $x_n$ if,
      for large values of $n$, $x_n$ is necessarily arbitrarily close to
      $L$. We then say that $L$ is the limit of the sequence of terms $x_n$
      and write $L = \lim \{x_n\}$.
\item The \emph{circumference} $C$ of a circle is the limit of the
      perimeters $p_n$ of regular inscribed polygons. That is,
      $C = \lim \{p_n\}$.
\item The number $C/d$, which is the same for all circles, will be denoted
      by the Greek letter $\pi$ (pi).
\item The \emph{area $S$ of a circle} is the limit of the areas of inscribed
      regular polygons. That is, $S = \lim \{S_n\}$.
\item Given a circular arc, let $l_n$ be the length of a regular polygonal
      path of $2^n$ sides inscribed in the arc. The \emph{length $l$ of the
      arc} is defined to be the following limit:
      $l = \mathrm{limit}\, \{l_n\}$.
\item On a circle of radius $r$ an arc of length $\pi r/180$ is subtended by
      a central angle of \emph{one degree} ($1^\circ$).
\item If on a circle of radius $r$ a central angle subtends an arc of length
      $L$, then the \emph{measure of the angle in degrees} is
      $(L/\pi r) \cdot 180^\circ$.
\end{applist}

%========================================================= theorems
\APPhead{SELECTED THEOREMS}{Selected Theorems}

The theorems selected below have been chosen to help you remember how the
ideas have been developed and \emph{in what order}. There is no attempt at
completeness, and in some cases the statement is less precise than in the
text; when this is the case, you should try to fill in the details or refer
back to the text.

\begin{applist}
\item Two different lines intersect in at most one point.
\item Let $O$ be any point of a line $l$. The point $O$ separates the line
      into two sets, called the two sides of the point $O$ on the line,
      having the following properties: (1) If $P$ and $Q$ belong to the same
      side of $O$, then either $OPQ$ or $OQP$. (2) If $P$ and $Q$ belong to
      different sides of $O$, then $POQ$.
\item If $P$ and $Q$ are both in the interior of $\ang{O}$, then the segment
      $PQ$ is in the interior of $\ang{O}$.
\item If rays $r_1$ and $r_2$ form $\ang{O}$, a ray $r$ from $O$, other than
      $r_1$ or $r_2$ lies either completely in the interior or completely in
      the exterior of $\ang{O}$.
\item If $P$ is in the interior of $\ang{O}$ and $Q$ is in the exterior of
      $\ang{O}$, then the segment $PQ$ contains a point of the angle.
\item If $AB$ and $A'B'$ are any two segments, either $AB > A'B'$, or
      $AB \cong A'B'$, or $AB < A'B'$, and exactly one of these holds.
\item If $AB > CD$ and $CD \cong EF$, then $AB > EF$.
\item If $AB \cong A'B'$, then $\sg{AB} = \sg{A'B'}$.
\item If $\sg{AB} = \sg{A'B'}$, then $AB \cong A'B'$.
\item $AB > CD$ if and only if $\sg{AB} > \sg{CD}$.
\item The S.A.S. theorem for congruent triangles.
\item The A.S.A. theorem.
\item The S.S.S. theorem.
\item For any two angles $\ang{A}$ and $\ang{B}$ either
      $\ang{A} < \ang{B}$, or $\ang{A} \cong \ang{B}$, or
      $\ang{A} > \ang{B}$, and exactly one of these is true.
\item All right angles are congruent to one another.
\item An exterior angle of a triangle is greater than either opposite
      interior angle.
\item In $\tri{ABC}$, $AB > AC$ if and only if $\ang{C} > \ang{B}$.
\item In $\tri{ABC}$, $\sg{AB} + \sg{BC} > \sg{AC}$.
\item If $l$ is any line and $A$ any point not on $l$, then there exists at
      least one line on $A$ parallel to $l$.
\item If parallel lines are cut by a transversal, the alternate interior
      angles are congruent.
\item The sum of the angles of a triangle is a straight angle.
\item If a transversal is cut by three parallel lines in congruent segments,
      the same is true for every transversal of these lines.
\item If for $\tri{ABC}$ and $\tri{A'B'C'}$, $\ang{A} \cong \ang{A'}$,
      $\ang{B} \cong \ang{B'}$, and $\ang{C} \cong \ang{C'}$, then the
      triangles are similar.
\item If $\sg{AB}/\sg{A'B'} = \sg{AC}/\sg{A'C'} = \sg{BC}/\sg{B'C'}$, then
      $\tri{ABC}$ is similar to $\tri{A'B'C'}$.
\item In a right triangle the altitude on the hypotenuse forms two right
      triangles similar to the given triangle.
\item In $\tri{ABC}$ with $\ang{C}$ a right angle,
      $\sg{AC}^2 + \sg{BC}^2 = \sg{AB}^2$.
\item The area of a triangle is one half the base times the altitude.
\item Suppose that circles with centers $A$ and $B$ each have radius
      $r > \frac{1}{2}\sg{AB}$. Then these circles intersect in two points
      on opposite sides of the line $AB$.
\item Exactly one circle is circumscribed about a regular polygon.
\item If regular polygons of $n$ sides are inscribed in two circles, then
      the perimeters of the polygons are proportional to the radii of the
      circles.
\item For a given circle let $p_n$ be the perimeter of a regular polygon of
      $n$ sides inscribed in the circle. The numbers $p_n$,
      $n = 3, 4, 5, \ldots$, form a sequence which has a limit.
\item The ratio $C/d$ (circumference to diameter) is the same for all
      circles.
\item If $S_n$ is the area of a regular polygon of $n$ sides inscribed in a
      circle, then $\mathrm{limit}\, \{S_n\}$ exists.
\item The area of a circle is equal to one half the circumference times the
      radius.
\item If two central angles are congruent, they subtend arcs of equal
      length.
\item If $\widehat{AB}$ and $\widehat{BC}$ are adjacent arcs (that is, if
      they have only point $B$ in common), then length $\widehat{ABC}$
      ${}={}$ length $\widehat{AB}$ ${}+{}$ length $\widehat{BC}$.
\item The measure of an angle inscribed in a circle is one half the measure
      of the central angle which subtends the same arc.
\end{applist}

\end{document}
